% SPDX-License-Identifier: GPL-3.0-or-later
% Copyright (C) 2026 Alessandro Greco
\chapter{观察流量：Wireshark 与日志}\label{cap:20-traffico}
\aggiornatoal{2026 年 9 月 23 日}{tor 0.4.9.13；tshark（Debian trixie）；实验网络}

\begin{obiettivi}
学完本章后，你将能够：
\begin{itemize}[nosep]
  \item 使用 \cmd{tcpdump} 捕获 Tor 流量，并借助 \cmd{tshark} 或 Wireshark 分析；
  \item 区分外部观察者能够看到的信息与仍受保护的信息；
  \item 利用会话密钥解密\emph{自己建立的}通道上的 TLS，并检查其中的信元（实验~12）；
  \item 通过 tor 的日志，逐步跟踪协议的运行过程。
\end{itemize}
\end{obiettivi}

前面几章，我们通过编写直接使用协议的代码，从\enquote{内部}研究协议。本章则从\enquote{外部}观察它，如同一个监听网络流量的人一样。这项练习有两个方向相反的目的：了解攻击者能从流量中推断出什么，以及获得一种独立验证自己实现的方法。

\begin{attenzione}[只观察自己的流量，而且只在实验环境中]
在意大利以及大多数其他国家，未经授权捕获他人的流量都可能违法（\cref{cap:25-legale}），同时也违背匿名研究的伦理原则（\cref{sec:rp-perche}）。本章所有抓包操作，都只针对本书的私有网络或我们自己建立的连接。
\end{attenzione}

% ---------------------------------------------------------------------------
\section{观察者能够看到什么？}\label{sec:traffico-osservatore}

观察客户端与入口守卫之间连接的人，例如互联网接入服务商，能够看到：
\begin{itemize}
  \item IP 地址和端口：客户端连接到某个地址；如果它不是网桥，该地址通常会出现在公开共识中，因此可以被识别为 Tor 中继；
  \item TLS 握手：在 TLS~1.3 下，证书经过加密传输，但客户端最初发送的消息仍有部分特征可见；
  \item TLS 记录的大小和时间特征。
\end{itemize}
观察者无法看到 cell 的内容、电路信息或目的地。然而，\emph{网站指纹识别}和流量关联攻击正是依靠数据大小与时间特征（\cref{cap:26-attacchi}）。内容被加密，并不意味着流量完全不泄露信息。

在 \cref{sec:lab-censura} 中，我们已经比较了能够被识别为 TLS 的普通 Tor 流量，以及不具有明显结构的 obfs4 流量。现在再进一步，尝试查看 TLS 内部。

% ---------------------------------------------------------------------------
\section{使用的工具}\label{sec:traffico-strumenti}

\begin{description}
  \item[\cmd{tcpdump}] 从网络接口捕获数据包，并保存到 \texttt{pcap} 文件。例如，在实验容器内，捕获发往 \texttt{test004r} ORPort 的流量：
\begin{lstlisting}[language=bash,numbers=none]
$ docker exec -u 0 torbook-network tcpdump -i lo \
    -w /lab/net/tls.pcap "tcp port 5104"
\end{lstlisting}
  \item[Wireshark 和 \cmd{tshark}] 分析抓包文件：前者提供图形界面，后者通过命令行使用。它们能够识别 TLS，但没有官方的 Tor 协议解析器。即使解密 TLS，Tor 信元仍然显示为原始字节。我们将使用 \file{labs/zh-Hans/代码/torlab/cells.py} 模块来解析这些字节。
\end{description}

\begin{figure}[htbp]
\centering
% SPDX-License-Identifier: GPL-3.0-or-later
% Copyright (C) 2026 Alessandro Greco
%% EPUB 无障碍替代文本。
\altfig{TCP 和 TLS 上的 Tor 封装示意图。外部观察者能看到 IP 地址、端口、流量大小、时间特征和 TLS 记录长度；没有 TLS 会话密钥则无法解析内部的 Tor cell。一个 TLS 记录可含多条 cell，也可能将一条 cell 分割到多个记录中，因此两者边界未必对齐。链路协议 v4 或更高版本中，固定长度 cell 为 514 字节，其中 5 字节为首部、509 字节为主体；还有可变长度 cell。只有 RELAY cell 的主体采用电路密钥逐层加密，PADDING 和 VPADDING 等 cell 不是 RELAY 消息。获得 TLS 密钥只能解析 cell，不能自动解密其 RELAY 主体。}
% 封装示意图：TLS 记录边界与 Tor cell 边界不必一致。
\begin{tikzpicture}[font=\sffamily\scriptsize,
  liv/.style={draw,rounded corners=2pt,minimum height=0.85cm,align=center},
  cel/.style={draw,fill=torviola!12,minimum height=0.7cm,minimum width=1.5cm,font=\sffamily\tiny}]
  \node[liv,fill=torgrigio!10,minimum width=10.6cm] (tcp) at (0,2.6) {TCP 字节流：分段长度可变};
  \node[liv,fill=torverde!12,minimum width=9.4cm] (tls) at (0,1.35) {TLS：记录与 cell 边界不对齐};
  \foreach \i/\x in {1/-3.6,2/-1.8,3/0,4/1.8,5/3.6} \node[cel] (c\i) at (\x,0) {514\,B};
  \node[anchor=west,font=\sffamily\tiny,text=torgrigio] at (4.6,0) {定长 v4+};
  \node[liv,fill=avviso!10,minimum width=4.6cm] (rel) at (-1.7,-1.4) {RELAY：509\,B 主体\\逐层加密};
  \node[liv,fill=torgrigio!8,minimum width=4.6cm] (pad) at (3.3,-1.4) {PADDING / VPADDING\\不是 RELAY};
  \draw[torgrigio] (c1.south west) -- (rel.north west); \draw[torgrigio] (c3.south east) -- (rel.north east);
  \node[anchor=west,align=left,font=\sffamily\tiny,text=torgrigio] at (5.8,2.6) {IP、端口、流量\\与时间可观察};
  \node[anchor=west,align=left,font=\sffamily\tiny,text=torgrigio] at (5.8,1.35) {TLS 记录长度\\在链路上可见};
  \node[anchor=west,align=left,font=\sffamily\tiny,text=torgrigio] at (5.8,0) {解密 TLS 后\\才能解析 cell};
  \node[anchor=west,align=left,font=\sffamily\tiny,text=torgrigio] at (5.8,-1.4) {RELAY 主体还需\\电路密钥};
  \node[font=\sffamily\tiny,text=torgrigio,align=center] at (0,-2.55) {仅为示意：TLS 记录与 cell 的边界不一定对齐。};
\end{tikzpicture}

\caption{Tor 通道示意：TCP 分段承载 TLS 记录，记录边界不一定与 cell 边界对齐。IP 地址、端口、长度和时间特征均可观察；解密 TLS 后才能读取 cell。链路协议 v4+ 的固定长度 cell 为 514 字节（首部 5 字节、主体 509 字节），但也存在可变长度 cell。只有 \texttt{RELAY} cell 主体由电路密钥逐层加密；填充 cell 并不是 RELAY 消息。}
\label{fig:incapsulamento}
\end{figure}

\subsection{解密自己建立的 TLS 连接}

TLS 使用外部观察者不具备的会话密钥。但是，\emph{直接参与}连接的一方可以导出这些密钥。浏览器和启用密钥日志功能的 Python 应用可以将会话密钥写入文件。本书的 \file{channel.py} 会在设置 \texttt{SSLKEYLOGFILE} 时显式配置 \texttt{SSLContext.keylog\_filename}，从而导出密钥。Wireshark 和 \cmd{tshark} 可以使用该文件解密抓包数据。实验~1 的 \file{channel.py} 遵循这一约定：运行前设置好环境变量即可。文件中可能包含以下内容：
\begin{lstlisting}[numbers=none]
# TLS secrets log file, generated by OpenSSL / Python
SERVER_HANDSHAKE_TRAFFIC_SECRET 4878224d3f9fb9aedb96d... 0da89dff6feb...
\end{lstlisting}
这些记录包含客户端随机数，以及对应 TLS~1.3 会话中各个阶段的秘密值。

\begin{approfondimento}[为什么这不构成对 Tor 的攻击？]
解密自己通道上的 TLS，不会揭示客户端原本不知道的信息。即使获得 TLS 会话密钥，\cella{RELAY} 信元仍然使用每一跳的独立密钥进行多层加密（\cref{cap:10-cifratura-relay}）。要读取它们，还需要只有客户端和对应中继才持有的电路密钥。TLS 只是洋葱的最外层，而不是唯一一层。
\end{approfondimento}

% ---------------------------------------------------------------------------
\section{实验：TLS 内部的信元}\label{sec:lab-traffico}

\begin{laboratorio}[12]{捕获、解密并解析一条通道}
\textbf{目标。}捕获我们自己的客户端与中继建立的通道，使用会话密钥解密，并通过 \file{labs/zh-Hans/代码/torlab/traffic.py} 还原信元序列。\\
\textbf{前置条件。}实验~1--3；实验网络已经运行；容器中已安装 \cmd{tcpdump} 和 \cmd{tshark}。\\
\textbf{预计时间。}1--2 小时。
\end{laboratorio}

执行实验~12 时，请先启动私有实验网络，并在仓库根目录打开两个终端。下面的抓包命令应在\emph{宿主机}上运行：它会在\emph{实验容器内部}以 root 身份执行 \cmd{tcpdump}，因为 Makefile 启动的普通容器 shell 使用非特权用户。第二个终端运行客户端时，保持抓包运行，最后使用 Ctrl+C 结束。捕获文件写入共享卷 \file{labs/zh-Hans/实验室/net/}。

\paragraph{步骤 1：抓包。}先在容器中启动抓包，然后从主机向 \texttt{test004r} 打开一条通道，在导出 TLS 会话密钥的同时建立一跳电路：
\begin{lstlisting}[language=bash,numbers=none]
$ PYTHONPATH=labs/zh-Hans/代码 SSLKEYLOGFILE=labs/zh-Hans/实验室/net/tls-keys.log python3 -c "
from torlab.laboratory import relay
from torlab.circuit import build
build([relay('004r')], cc=False)"
\end{lstlisting}
前述命令已在同一行指定 \texttt{PYTHONPATH}，因此无需另外设置导入路径。

请在采用主机网络的 Linux 系统上，从\emph{仓库根目录}执行 Python 命令。\texttt{PYTHONPATH} 使 \file{labs/zh-Hans/代码/} 下的 \file{torlab} 包可以导入；\file{labs/zh-Hans/代码/torlab/channel.py} 会显式读取 \texttt{SSLKEYLOGFILE} 并开启 TLS 会话密钥导出。密钥日志包含敏感会话秘密，不得公开，并应在实验后删除。

\paragraph{步骤 2：识别自己的连接。}实验网络中的其他中继也会连接同一个端口，但它们的会话密钥不同。通过 \cmd{tshark} 尝试解密并逐条跟踪 TCP 流，找出能够成功解密的连接：
\begin{lstlisting}[language=bash,numbers=none]
$ docker exec torbook-network tshark -r /lab/net/tls.pcap \
    -o tls.keylog_file:/lab/net/tls-keys.log \
    -q -z follow,tls,raw,2
\end{lstlisting}
TCP 流编号 \texttt{2} 仅来自以前的一份示例抓包，并非固定编号。可利用 \cmd{tshark} 的 \texttt{-T fields -e tcp.stream} 查看候选流，再选择能够成功解密的流。下面的输出是历史示例，证书长度和电路 ID 会随着运行而改变。
输出会以十六进制给出两个方向上的解密数据，其中一个方向通过制表符缩进。

\paragraph{步骤 3：检查 cell。}\texttt{cells\_from\_follow} 函数会重建两个方向上的字节流，并把它们交给实验~2 的 \texttt{decode}。每个方向都有自己的协议版本协商：第一个信元是带有两字节 CircID 的 \cella{VERSIONS}。只有在版本协商完成后，CircID 才使用四字节格式。结果如下：
\begin{lstlisting}[numbers=none]
client -> relay  CircID=0x0000 VERSIONS 消息体=4 字节
relay  -> client CircID=0x0000 VERSIONS 消息体=6 字节
relay  -> client CircID=0x00000000 CERTS 消息体=1443 字节
relay  -> client CircID=0x00000000 AUTH_CHALLENGE 消息体=36 字节
relay  -> client CircID=0x00000000 NETINFO 消息体=509 字节
client -> relay  CircID=0x00000000 NETINFO 消息体=509 字节
client -> relay  CircID=0xc244f1f6 CREATE2 消息体=509 字节
relay  -> client CircID=0xc244f1f6 CREATED2 消息体=509 字节
\end{lstlisting}
从外部观察到的序列，恰好就是 \cref{fig:sessione} 和 \cref{cap:07-canali} 所描述的过程。客户端的 \cella{VERSIONS} 信元消息体长四字节，因为 \file{channel.py} 只提供版本 4 和 5。电路 ID \texttt{0xc244f1f6} 的最高有效位为 1，因为通道由客户端发起（\cref{sec:circid}）。

\begin{attenzione}[初版实验中的两个错误]
最初的分析犯过两个错误，直到把输出与预期序列比较后才发现。首先，我们错误地把同一个协商版本用于两个方向，因而误解了中继的 \cella{VERSIONS} 信元。其次，我们靠猜测判断 \cmd{tshark} 对哪个方向做了缩进，结果将数据块归入了错误方向。正确的判断规则很简单：第一个数据块始终是客户端发送的 \cella{VERSIONS} 信元。
\end{attenzione}

% ---------------------------------------------------------------------------
\section{tor 日志}\label{sec:traffico-log}

抓包能展示网络中\emph{传输了什么}，而 tor 的日志则帮助解释\emph{为什么}。在 \texttt{info} 级别，tor 会记录电路建立阶段、守卫选择、描述符下载和发布等活动；在 \texttt{debug} 级别，记录还可以细化到单个 cell。我们已经利用日志确认客户端查询了哪个 HSDir（\cref{sec:lab-onion}），以及洋葱服务何时更新描述符。

下面是几点实用建议：
\begin{itemize}
  \item \texttt{debug} 会产生大量文本。可只为特定模块启用，例如使用 \texttt{Log [circ]debug file circ.log}，把详细记录限制在电路事件。
  \item 在真实网络中，\opt{SafeLogging} 会掩盖地址信息。客户端日志不应记录用户访问的目的地。
  \item 公开中继的日志可能包含其他人的个人信息，并产生 \cref{cap:25-legale} 讨论的法律后果；只应保留严格必要的数据。
\end{itemize}

% ---------------------------------------------------------------------------
\section*{本章小结}
\begin{itemize}
  \item 连接观察者可以看到 IP 地址、TLS 握手、数据包大小和时间特征，但无法直接看到 Tor 信元或目的地。
  \item \texttt{SSLKEYLOGFILE} 可以让我们解密自己建立的通道上的 TLS，从而查看 cell；不过 \cella{RELAY} 信元仍受多层加密保护。
  \item 分析抓包数据为验证协议实现提供了独立方法。
  \item tor 的 \texttt{info} 和 \texttt{debug} 日志可解释程序所做的选择；真实使用时应尽量减少日志，并妥善保护。
\end{itemize}

% ---------------------------------------------------------------------------
\section*{练习}

\begin{esercizio}[1]\label{es:20-dimensioni}
在一次 Tor 下载流量抓包中，TLS 记录的长度并不全是 514 字节的整数倍。既然 cell 长度固定，为什么会出现这种情况？
\end{esercizio}

\begin{esercizio}[2]\label{es:20-relay}
扩展实验~12，使用 \file{circuit.py} 提供的密钥解密自己电路上的 \cella{RELAY} 信元。构建三跳电路的过程中，能够观察到哪些 relay 消息？
\end{esercizio}

\begin{esercizio}[2]\label{es:20-padding}
在不产生应用程序流量的情况下，捕获实验客户端通道上持续一分钟的通信。统计两个方向的 TLS 记录数量，并将观察结果与 \cref{cap:07-canali} 中的通道填充机制联系起来。
\end{esercizio}

\begin{esercizio}[3]\label{es:20-dissector}
使用 Lua 编写一个简单的 Wireshark 协议解析器，识别解密后数据中的 Tor 信元，并显示 CircID 和命令类型。
\end{esercizio}

% ---------------------------------------------------------------------------
\section*{延伸阅读}
\begin{itemize}
  \item Wireshark 文档中有关利用会话密钥日志文件解密 TLS 的说明。
  \item tor 手册的 \opt{Log} 条目，介绍日志模块和详细程度。
\end{itemize}
